\mypart{Напоминания о~булевой алгебре}

\section{Высказывания и~булева алгебра}

Говоря неформально, высказывание~"--- это что-то, на что можно ответить «да» или «нет».

Формально говоря,
\begin{definition}
	\tdef{Высказывание}~"--- это элемент~$\set*{0, 1}$.
\end{definition}

Элемент~$0$ трактуем как ложь, а~$1$~"--- как истину.

Иногда нам будет удобно использовать отображение из какого-либо множества в~$\set*{0, 1}$.

%\begin{definition}
%	\tdef{Булева алгебра}~"--- это множество $\set*{0, 1}$ с~заданными на нём операциями, переводящими $\set*{0, 1}$ в~$\set*{0, 1}$.
%	\label{<+label+>}
%\end{definition}<++>

Напомним различные операции в~булевой алгебре.

Унарных операций (то есть таких, которые принимают всего один аргумент) две~"--- тождественная и~инверсия:
\begin{align*}
	\begin{array}{c|c}
		x & \id(x)
		\\\hline
		0 & 0
		\\
		1 & 1
	\end{array}
	&&
	\begin{array}{c|c}
		x & \overline{x}
		\\\hline
		0 & 1
		\\
		1 & 0
	\end{array}
\end{align*}

Если $x$ и~$y$~"--- различные высказывания, то
\[
	\begin{array}{*{3}{>{$\displaystyle}p{0.20\linewidth}<{$}}}
		\begin{array}{cc|c}
			x & y & x \wedge y
			\\\hline
			0 & 0 & 0
			\\
			0 & 1 & 0
			\\
			1 & 0 & 0
			\\
			1 & 1 & 1
		\end{array}
		&
		\begin{array}{cc|c}
			x & y & x \vee y
			\\\hline
			0 & 0 & 0
			\\
			0 & 1 & 1
			\\
			1 & 0 & 1
			\\
			1 & 1 & 1
		\end{array}
		&
		\begin{array}{cc|c}
			x & y & x \implies y
			\\\hline
			0 & 0 & 1
			\\
			0 & 1 & 1
			\\
			1 & 0 & 0
			\\
			1 & 1 & 1
		\end{array}
		\\[10ex]
		\begin{array}{cc|c}
			x & y & x \oplus y
			\\\hline
			0 & 0 & 0
			\\
			0 & 1 & 1
			\\
			1 & 0 & 1
			\\
			1 & 1 & 0
		\end{array}
		&
		\begin{array}{cc|c}
			x & y & x \equiv y
			\\\hline
			0 & 0 & 1
			\\
			0 & 1 & 0
			\\
			1 & 0 & 0
			\\
			1 & 1 & 1
		\end{array}
	\end{array}
\]

В~логике есть ещё и~импликация; она соответствует выражению <<если-то>>: $P \implies Q$ означает <<если истинно~$P$, то истинно~$Q$>>.

При этом выражение $P \implies Q$ можно трактовать ещё так:
\begin{itemize}
	\item <<чтобы выполнялось~$P$, необходимо, чтобы выполнялось~$Q$>>;
	\item <<для выполнения~$Q$ достаточно выполнения~$P$>>.
\end{itemize}

\begin{note}
	Напомню, что импликации $(P \implies Q)$ и~$(\overline{Q} \implies \overline{P})$ эквивалентны.
\end{note}

\begin{note}
	Операция~<<$\equiv$>> называется \tdef{эквивалентностью}; также для неё используют знак <<$\Longleftrightarrow$>> и~читают <<тогда и~только тогда, когда>>.
\end{note}


\section{Кванторы}

Если выражение~$P(x)$, зависящее от высказывания~$x$, выполняется при всех~$x$, то пишут
\[
	\forall x
	\quad
	P(x)
	,
\]
или
\[
	\forall x
	\quad
	\bigl(
		P(x)
	\bigr)
	.
\]
Я~использую для разделения кванторов и~высказываний значок~$\vdash$, который пришёл из логики:
\begin{equation}
	\label{eq:Logic:forall}
	\forall x \holds P(x)
	.
\end{equation}
Это читается так: <<для любого $x$ выполняется~$P(x)$>>.

Если выражение~$P(x)$ выполняется хотя бы при одном~$x$, то пишут
\[
	\exists x 
	\quad
	P(x),
\]
или
\[
	\exists x
	\quad
	\bigl(
		P(x)
	\bigr)
	,
\]
или
\begin{equation}
	\label{eq:Logic:exists}
	\exists x
	\suchthat
	P(x)
	.
\end{equation}
Это читается так: <<существует элемент~$x$, такой, что выполняется~$P(x)$>>.
Обрати внимание, что таких элементов, при которых выполняется~$P$, может быть
несколько~"--- в~том числе все.

Значок <<$\forall$>> читается как <<для любого>> или <<любой>>; значок
<<$\exists$>>~"--- как <<существует>>. Эти обозначения получаются из первых букв
слов <<All>> и~<<Exists>> отражением их относительно левой нижней точки буквы.
Значок <<$\vdash$>> можно читать как <<такой, что>>.

Вместе со значком существования~$\exists$ в~некоторых случаях используется знак
<<$!$>>~"--- <<единственный>>:
\[
	\exists! x \suchthat P(x)
\]
<<существует \emph{единственный} $x$, такой, что выполняется~$P(x)$>>.

\begin{note}
	В~выражении~\eqref{eq:Logic:exists} <<$\exists x \suchthat P(x)$>>
	(<<существует~$x$, такой, что выполнено~$P(x)$>>) подразумевается, что~$P(x)$
	может быть выполнено при более чем одном значении~$x$~"--- в~том числе и~при
	всех. Выражение~\eqref{eq:Logic:exists} гарантирует только, что существует
	\emph{хотя бы один} элемент~$x$, для которого выполняется~$P(x)$.
\end{note}

\begin{note}
	Порядок кванторов имеет значение. Если встречается квантор существования, то
	переменная, к~которой он прицеплен, зависит от всех переменных, встретившихся
	до неё. К~примеру, в~выражении
	\[
		\forall x\ 
		\exists y\ 
		\forall z\ 
		\exists v
	\]
	\begin{itemize}
		\item переменная~$x$ не зависит ни от чего;
		\item переменная~$y$ зависит от~$x$;
		\item переменная~$z$ не зависит ни от чего;
		\item переменная~$v$ зависит от~$x$, $y$, $z$.
	\end{itemize}
\end{note}

\begin{note}
	В~жизни встречается знак равенства по определению~"--- <<$:=$>> или <<$=:$>>.
	Двоеточие ставится со стороны определяемого объекта.
\end{note}

\section{Отрицание}

\begin{definition}
	\label{def:neg_quantors_simple}
	Определим отрицание выражений с~кванторами:
	\begin{align*}
		\neg
		\bigl(
			\forall x
			\holds
			P(x)
		\bigr)
		&:=
		\bigl(
			\exists x
			\suchthat
			\neg
			P(x)
		\bigr)
		\refstepcounter{equation}
		\tag{\theequation.1}
		\label{eq:sets:forall:x}
		,
		\\
		\neg
		\bigl(
			\exists x
			\suchthat
			P(x)
		\bigr)
		&:=
		\bigl(
			\forall x
			\holds
			\neg
			P(x)
		\bigr)
		.
		\tag{\theequation.2}
		\label{eq:sets:forall:x}
	\end{align*}
\end{definition}
